\section{Mettre les choix de sécurité en regard des risques OWASP}

La sécurité de FeastVerse n'a pas été conçue à partir d'une checklist OWASP appliquée formellement pendant le développement. Il serait donc excessif de présenter le projet comme \og conforme OWASP \fg{}. En revanche, plusieurs mécanismes effectivement implémentés répondent directement à des familles de risques identifiées par l'OWASP Top 10 \parencite{owasp-top10}. Cette mise en correspondance permet de mieux expliciter la portée des choix techniques déjà décrits dans le chapitre.

Le premier point concerne le contrôle d'accès. L'application ne se limite pas à vérifier qu'un utilisateur est authentifié : les rôles, la propriété des ressources et certaines règles métier déterminent les actions autorisées. Les annotations Spring Security protègent les opérations réservées aux modérateurs ou administrateurs, tandis que les services vérifient les droits liés au propriétaire lorsque le rôle seul ne suffit pas. Cette organisation contribue à réduire le risque de contrôle d'accès défaillant, sans prétendre supprimer à elle seule toutes les possibilités de contournement.

L'authentification par \terme{jwt} est également intégrée à une chaîne de validation complète : extraction du jeton, vérification de sa validité, rechargement de l'utilisateur et initialisation du contexte de sécurité. Les contrats HTTP sont par ailleurs séparés du modèle de persistance au moyen de \terme{dto}, afin de limiter les informations exposées et de contrôler les données reçues. La validation des entrées et la gestion centralisée des erreurs complètent cette démarche en évitant qu'une requête invalide soit traitée comme un cas nominal.

Enfin, la séparation entre contrôleurs, services et repositories limite la dispersion des règles de sécurité dans le code. Elle facilite la revue des autorisations et la localisation des décisions métier. Cette structure ne remplace toutefois pas des contrôles de sécurité dédiés. Le projet ne comprend pas, dans son état actuel, de campagne automatisée d'analyse de vulnérabilités, de test d'intrusion, de rotation industrialisée des secrets ou de revue complète des dépendances selon une politique formalisée.

Cette lecture par rapport à OWASP permet donc surtout de distinguer deux niveaux : d'une part des pratiques de sécurisation réellement présentes dans FeastVerse ; d'autre part les vérifications complémentaires qui seraient nécessaires avant de présenter la solution comme prête pour une exploitation exposée à des utilisateurs réels. Elle renforce ainsi la justification de la compétence C2.4 sans transformer une référence de bonnes pratiques en certification de sécurité.